\section{Geophysical implications}

\subsection{Equivalent layer vs. physical nature of the observed data}

The general expression proposed by \cite{oliveirajr_etal2026_1} to define $p_{\varepsilon}$ in 
the Fourier domain (Eq. \ref{eq:p-Fourier}) allows setting different classes of 
equivalent layers.
In the particular case where $\alpha_{a} = \alpha_{d}$, for example, we obtain a Neumann's 
equivalent layer, which preserves the physical relationship between the observed data and the 
true sources. It means that, if the observed data correspond to the $y$-component of the 
magnetic induction field of the true sources, 
then the potential field predicted by the equivalent layer is also the $y$-component of the 
magnetic induction field of the equivalent sources. Most of the equivalent-layer methods use 
such layers to perform potential-field transformations.
On the other hand, layers $p_{\varepsilon}$ obtained with $\alpha_{a} \ne \alpha_{d}$ do not
preserve the physical relationship between the observed data and the true sources. 
This is the general case, for example, of layers obtained with $\alpha_{a} = k$, which are 
called Dirichlet's equivalent layers \citep{oliveirajr_etal2026_1}. Such layers are commonly used 
to perform the continuation of potential-field data. Obviously that Dirichlet's equivalent 
layers coincide with Neumann's equivalent layers for the particular case in which 
$\alpha_{d} = k$. 

One important aspect of Eq. \ref{eq:p-Fourier} is that it provides the theoretical 
foundation for common practices that have been empirically proved by several studies: 
equivalent layers are independent of the physical nature of the observed potential-field data 
and can be conveniently chosen according to the characteristics of the problem. Many authors 
have proposed equivalent-layer methods based on this empirical principle
\citep[e.g.,][]{cordell1992, guspi_novara2009, li_etal2014, soler_uieda2021}, which is 
theoretically justified here for the first time by using the framework proposed by
\cite{oliveirajr_etal2026_1}.

\subsection{Conversion of gravity disturbance into gravity-gradient data}

The common two-step approach to convert 
gravity disturbance data into the gravity-gradient tensor components
\citep[e.g.,][]{piauilino_etal2019_seg, oliveirajr_etal2023, piauilino_etal2025b} via equivalent-layer 
technique empirically presumes the existence of a single layer of 
monopoles that satisfactorily reproduces the gravitational potential and 
its derivatives of first and second orders, which are represented by the 
gravity disturbance and gravity-gradient data, respectively. 
Specifically, the first step in this approach consists in solving Eq. 
\ref{eq:predicted-data-continuous-layer} 
for $p_{\varepsilon}$ by using $D_{a} = \nicefrac{\partial}{\partial z}$, 
so that the predicted potential field $d_{\varepsilon}$ approximates the observed potential 
field $d^{o}$ (Eq. \ref{eq:observed-data-space}).
Then, in the second step, the estimated $p_{\varepsilon}$ is used to compute each 
gravity-gradient component with
Eq. \ref{eq:transformed-field-continuous-layer}, where $D_{t}$ is the corresponding 
differential operator associated with the gravity-gradient component.
For example, if the desired component is the $xy$, so 
$D_{t} = \nicefrac{\partial^{2}}{\partial x \, \partial y}$.

To understand why this approach works, let us first note that 
the differential operator associated with the observed gravity disturbance data (Eq. 
\ref{eq:observed-data-space}) is given by $D_{d} \equiv D(1, \hat{\mathbf{z}}) = 
\tfrac{\partial}{\partial z}$ in the space domain
or $\alpha_{d} \equiv \alpha(1, \hat{\mathbf{z}}) = k$ in the Fourier domain (Eq. 
\ref{eq:diff-operators-generic}),
where $\hat{\mathbf{z}} = (0, 0, 1)^{\top}$.
The first step consists in using $D_{a} = D_{d}$ in Eq. 
\ref{eq:predicted-data-continuous-layer} and 
estimating $p_{\varepsilon}$ so that the corresponding predicted potential field 
$d_{\varepsilon}$ approximates
the observed potential field $d^{o}$ (Eq. \ref{eq:observed-data-space}).
In this case, we conclude from equations \ref{eq:diff-operators-generic} and 
\ref{eq:p-Fourier} that $p_{\varepsilon}$ 
approximates the vertical derivative $\tfrac{\partial U}{\partial z}$ of $U$
(Eq. \ref{eq:potential-U-space}) on the surface $\mathcal{S}''$.

The second step to convert gravity disturbance data into gravity-gradient data involves using 
$p_{\varepsilon}$ to compute the gravity-gradient components, which are represented by the 
potential field $f$, with specific operators $D_{f}$ (Eq. 
\ref{eq:transformed-field-space}). 
Since our $p_{\varepsilon}$ (Eq. \ref{eq:p-Fourier}) is defined with $\alpha_{a} = 
\alpha_{d}$, it is a Neumann's equivalent layer \citep{oliveirajr_etal2026_1} and, consequently, 
any transformed potential field $f_{\varepsilon}$ 
(Eq. \ref{eq:transformed-field-continuous-layer}) must be defined with $D_{t} = D_{f}$, 
where $D_{f}$ is the space-domain 
operator associated with $f$ (Eq. \ref{eq:transformed-field-space}).
Consider, for example, that the desired potential field $f$ (Eq. 
\ref{eq:transformed-field-space}) is the gradient 
component $xy$. This component is defined by  
$D_{f} \equiv D(1, \hat{\mathbf{x}}) \, D(1, \hat{\mathbf{y}}) = \tfrac{\partial^{2}}{\partial x \, \partial y}$ in the
space domain, or 
$\alpha_{f} \equiv \alpha(1, \hat{\mathbf{x}}) \, \alpha(1, \hat{\mathbf{y}}) = \mathrm{i}^{2} \, k_{x} \, k_{y}$
in the Fourier domain (Eq. \ref{eq:diff-operators-generic}), where 
$\hat{\mathbf{x}} = (1, 0, 0)^{\top}$ and
$\hat{\mathbf{y}} = (0, 1, 0)^{\top}$.
In this case, the predicted $xy$ gradient component is obtained by evaluating 
$f_{\varepsilon}$ (Eq. \ref{eq:transformed-field-continuous-layer-Fourier}) with 
$D_{t} = \tfrac{\partial^{2}}{\partial x \, \partial y}$.
The remaining gradient components are evaluated in the same way, by using their corresponding 
operators $D_{f}$.

Thus, the common two-step approach to convert gravity disturbance into gravity
gradient data via equivalent-layer technique 
\citep[e.g.,][]{piauilino_etal2019_seg, oliveirajr_etal2023, piauilino_etal2025b} 
can be theoretically explained as an application of a particular Neumann's equivalent layer.

\subsection{Joint processing of gravity-gradient data}

Similarly to the previous approach to convert gravity disturbance data into gravity gradient 
data, the common approach to joint processing gravity gradient data 
\citep[e.g.,][]{boggs-dransfield2004,barnes_lumley2011, martinez_li_2016} via equivalent-layer 
technique empirically presumes that a single equivalent layer $p_{\varepsilon}$ can be used to 
predict all gravity gradient components, i.e., that
\begin{equation}
	d^{uv}_{\varepsilon}(\mathbf{r}) = 
	\int\limits_{-\infty}^{\infty}\int\limits_{-\infty}^{\infty}
	p_{\varepsilon}(\mathbf{r}'') \, 
	\tfrac{\partial^{2} \Psi}{\partial u \, \partial v}(\mathbf{r}, \mathbf{r}'') \:\: 
	d\mathcal{S}'' \: , \quad z < z''_{0} \: ,
	\label{eq:gravity_gradient_layer}
\end{equation}
where $d^{uv}_{\varepsilon}$ approximates the gradient component $uv$, with 
$u, v \in \{x, y, z\}$.
This approach considers a set of equations defined by \ref{eq:gravity_gradient_layer}, each one
for a given gravity gradient component $uv$, and simultaneously solves then for the single 
layer $p_{\varepsilon}$.
Again, the existence of such layer can be deduced from Eq. \ref{eq:p-Fourier}, for the 
particular case in which 
$\alpha_{a} = \alpha_{d}$. In this case, $p_{\varepsilon}$ is a Neumann's equivalent layer 
\citep{oliveirajr_etal2026_1}.
Consider, for example, the gravity gradient component $xz$. In this case, 
the observed potential field is defined by 
$D_{d} \equiv D(1, \hat{\mathbf{x}}) \, D(1, \hat{\mathbf{z}}) = \tfrac{\partial^{2}}{\partial x \, \partial z}$ in the
space domain, or 
$\alpha_{d} \equiv \alpha(1, \hat{\mathbf{x}}) \, \alpha(1, \hat{\mathbf{z}}) = \mathrm{i} \, k_{x} \, k$
in the Fourier domain (Eq. \ref{eq:diff-operators-generic}), where 
$\hat{\mathbf{x}} = (1, 0, 0)^{\top}$ and
$\hat{\mathbf{y}} = (0, 0, 1)^{\top}$.
From Eq. \ref{eq:p-Fourier}, we conclude that the layer $p_{\varepsilon}$ obtained with 
$\alpha_{a} = \alpha_{d}$ represents the vertical derivative 
$\tfrac{\partial U_{\varepsilon}}{\partial z}$ of the auxiliary potential 
$U_{\varepsilon}$ (Eq(s). \ref{eq:potential-U-epsilon-Fourier-layer-1} and 
\ref{eq:potential-U-epsilon-Fourier-layer-2}) 
on the surface $\mathcal{S}''$.
The same conclusion is obtained for any other gravity gradient component defined with a 
different $uv$, which shows that there exists a single equivalent layer $p_{\varepsilon}$ that 
simultaneously fits all gravity gradient components.
We show here that the common approach to joint processing gravity gradient data 
\citep[e.g.,][]{boggs-dransfield2004,barnes_lumley2011, martinez_li_2016} can be explained in 
terms of a Neumann's equivalent layer.

\subsection{The existence of a planar dipole layer}

A predominant trend in previous studies utilizing a planar equivalent layer of dipoles for 
magnetic-field transformations 
is to treat the layer as a purely mathematical abstraction employed for processing total-field 
anomaly data, a perspective exemplified by the remarks of \cite{pedersen1989} concerning the 
research of \cite{bhattacharyya_chan1977}.
Specifically, Pedersen pointed out that the first fundamental existence theorem for the 
Dirichlet's problem does not postulate that a general solution can be written as a layer of 
dipoles magnetized along its 
normal direction. In essence, Pedersen underscored that a dipole layer lacks physical 
significance.
In response, Chan affirmed that such a layer was employed in his work 
solely to address an integral equation, i.e., as a purely mathematical solution.

Since its inception in the 1960s, the equivalent-layer technique in applied geophysics has 
garnered relatively scant attention in terms of exploring the theoretical attributes of planar 
dipole layers capable of faithfully reproducing total-field anomaly data. 
This inquiry usually aims to discern the physical significance and possible link of such 
layers with boundary problems of potential theory.

Here, we leverage the theoretical framework proposed by \cite{oliveirajr_etal2026_1} to infer 
within the purview of applied geophysics, for the first time, the existence of planar dipole 
layers capable of accurately and simultaneously reproducing the magnetic scalar potential and 
induction field components along any direction.
We first consider that the observed potential field $d^{o}$ (Eq. 
\ref{eq:observed-data-space}) is satisfactorily described by the differential operator 
$D_{d} \equiv D(1, \hat{\mathbf{u}}) \, D(1, \hat{\mathbf{v}})$
(Eq. \ref{eq:diff-operators-generic}), where $\hat{\mathbf{u}}$ defines the constant 
direction of the main geomagnetic field at the study area and $\hat{\mathbf{v}}$ the constant 
total-magnetization direction of the true magnetic sources.
From Eq. \ref{eq:diff-operators-generic}, we can show that 
(i) $D_{d} = \hat{\mathbf{u}}^{\top} \mathbf{H}_{U}(\mathbf{r}) \, \hat{\mathbf{v}}$,
where $\mathbf{H}_{U}$ is the Hessian matrix \citep[e.g.,][p. 158]{kaplan2002} of
the potential $U$ (Eq. \ref{eq:potential-U-space}) and (ii) its corresponding 
Fourier domain representation is given by 
$\alpha_{d} = \left(\mathbf{k}^{\top} \hat{\mathbf{u}} \right)\left(\mathbf{k}^{\top} \hat{\mathbf{v}} \right)$.

The common equivalent-layer approach to process total-field anomaly data consists in 
estimating an equivalent layer $p_{\varepsilon}$ by solving the Eq. 
\ref{eq:predicted-data-continuous-layer},
with $D_{a} = \hat{\mathbf{u}}^{\top} \mathbf{H}_{\Psi}(\mathbf{r}) \, \hat{\mathbf{h}}$, 
where $\mathbf{H}_{\Psi}$ is the Hessian matrix \citep[e.g.,][p. 158]{kaplan2002} of the 
function $\Psi$ (Eq. \ref{eq:Psi}) and $\hat{\mathbf{h}}$ is a constant unit vector not 
necessarily equal to $\hat{\mathbf{v}}$. From Eq. \ref{eq:diff-operators-generic}, we can 
show that the corresponding Fourier domain operator is given by 
$\alpha_{a} = \left(\mathbf{k}^{\top} \hat{\mathbf{u}} \right)\left(\mathbf{k}^{\top} \hat{\mathbf{h}} \right)$.
Then, by using the predefined $\alpha_{d}$ and $\alpha_{a}$ in Eq. \ref{eq:p-Fourier}, we 
obtain the following equivalent layer:
\begin{equation}
	\mathcal{F} [ p_{\varepsilon}(\mathbf{r}'') ] =
	k \, \tfrac{
		\left(\mathbf{k}^{\top} \hat{\mathbf{v}} \right)
		}{
		\left(\mathbf{k}^{\top} \hat{\mathbf{h}} \right)
		} \,
	\mathcal{F} \left[ U_{\varepsilon}(\mathbf{r}'') \right] \: .
	\label{eq:p-Fourier-dipoles}
\end{equation}

The particular equivalent layer given by Eq. \ref{eq:p-Fourier-dipoles} represents a 
planar layer of dipoles with constant magnetization direction $\hat{\mathbf{h}}$ and was 
deduced by considering
that $\hat{\mathbf{u}}$ is a unit vector having the same constant direction of the main 
geomagnetic field at the study area, i.e., by considering that the observed potential field is 
the approximated total-field anomaly \citep[e.g.,][p. 179]{blakely1996}.
Let us now consider that the observed potential field data is a Cartesian component of the 
magnetic induction field. In this case, we can simply replace $\hat{\mathbf{u}}$ by the unit 
vector along the direction of 
magnetic induction field component. For example, if the observed potential field is the $x$ 
component of the magnetic induction field, then we have
(i) $D_{d} = \hat{\mathbf{x}}^{\top} \mathbf{H}_{U}(\mathbf{r}) \, \hat{\mathbf{v}}$, 
(ii) $\alpha_{d} = \left(\mathbf{k}^{\top} \hat{\mathbf{x}} \right)\left(\mathbf{k}^{\top} \hat{\mathbf{v}} \right)$,
(iii) $D_{a} = \hat{\mathbf{x}}^{\top} \mathbf{H}_{U}(\mathbf{r}) \, \hat{\mathbf{h}}$ and
(iv) $\alpha_{a} = \left(\mathbf{k}^{\top} \hat{\mathbf{x}} \right)\left(\mathbf{k}^{\top} \hat{\mathbf{h}} \right)$,
where $\hat{\mathbf{x}} = (1, 0, 0)^{\top}$. For the cases in which the observed potential field is the $y$ or 
$z$ component of the magnetic induction field, then we replace $\hat{\mathbf{x}}$ by the 
unit vector $\hat{\mathbf{y}} = (0, 1, 0)^{\top}$ or $\hat{\mathbf{z}} = (0, 0, 1)^{\top}$,
respectively. For these three Cartesian components of the magnetic induction field, we conclude from Eq. 
\ref{eq:p-Fourier} that the resultant equivalent layer is that defined by Eq. \ref{eq:p-Fourier-dipoles}.
It can also be easily verified that the same $p_{\varepsilon}$ given by Eq. 
\ref{eq:p-Fourier-dipoles} can be equally deduced if we consider that the 
observed potential field is the magnetic scalar potential. 
In this case, we have $\alpha_{d} = \left(\mathbf{k}^{\top} \hat{\mathbf{v}} \right)$ and 
$\alpha_{a} = \left(\mathbf{k}^{\top} \hat{\mathbf{h}} \right)$. 

The main conclusion obtained here is that a planar equivalent layer of dipoles $p_{\varepsilon}$ defined according 
to Eq. \ref{eq:p-Fourier-dipoles} simultaneously 
reproduces the magnetic scalar potential and any magnetic induction field component (including the approximated total-field anomaly) 
due to sources with constant total-magnetization direction defined by $\hat{\mathbf{v}}$, even if the layer has a constant magnetization 
direction defined by a unit vector $\hat{\mathbf{h}} \ne \hat{\mathbf{v}}$. 
This theoretical result explains, for example, why we can use a planar equivalent layer of dipoles with arbitrary constant 
magnetization direction $\hat{\mathbf{h}}$ to compute the magnetic induction field components or the amplitude of the 
magnetic induction field due to the true sources without presuming any information about the constant direction of their 
total magnetization \cite[e.g,][]{li_etal_2021}.


\subsection{Computing the RTP by using a dipole layer}


According to the pioneer work of \cite{silva1986}, which is widely followed in applied geophysics, 
the RTP in the space domain is calculated from total-field anomaly data in two steps.
The first consists in estimating the magnetic-moment intensity distribution on a planar layer 
of dipoles (or doublets) with constant magnetization direction equal to that of the true sources.
The second step consists in computing the $z$-component of the magnetic induction field due 
to the estimated layer by setting its magnetization along the vertical direction.
It is usually verified a posteriori that this computed field approximates the true RTP anomaly.

We call attention to the fact that \cite{silva1986} cited the work of \cite{roy1962} to justify that ``the magnetic 
field due to arbitrary sources can be exactly fitted by the magnetic field due to a horizontal equivalent
layer of dipoles''. 
However, Roy neither discussed the existence of a dipole layer that fits the observed total-field 
anomaly data nor its use to retrieve the RTP anomaly, but 
presented a general outline to deduce the solutions of the Dirichlet's and Neumann's problems on a plane 
located above the sources.
In other words, the two-step approach proposed by \cite{silva1986}, which is widely accepted by the scientific community 
and we know to be correct from practice 
\citep[e.g.,][]{leao-silva_1989, guspi_novara2009, oliveirajr_etal2013, li_etal2014, reis_etal2020, takahashi-etal2022, piauilino_etal2025}, 
is an empirical generalization of \cite{roy1962}. 
It is also important to point out that the mathematical equivalence between the space-domain approach to compute the 
RTP via equivalent-layer technique and the classical approach directly using the Fourier transform of the observed 
total-field anomaly data \citep[e.g.,][p. 331]{blakely1996} have never been clearly established in the literature.
Here, we provide a general approach aiming at filling these theoretical gaps.

From a theoretical point of view, the two-step approach to compute the RTP in the space domain 
consists in first setting a planar layer of dipoles according to Eq. \ref{eq:p-Fourier-dipoles}, 
by considering the unit vector $\hat{\mathbf{h}}$ equal to $\hat{\mathbf{v}}$, which defines 
the constant total-magnetization direction of the true sources.
In this case, the ratio $\nicefrac{\left(\mathbf{k}^{\top} \hat{\mathbf{v}} \right)}{\left(\mathbf{k}^{\top} \hat{\mathbf{h}} \right)} = 1$ 
and we can easily verify that the $\mathcal{F} [ p_{\varepsilon}(\mathbf{r}'') ]$ given by 
Eq. \ref{eq:p-Fourier-dipoles} assumes the form of a Neumann's equivalent layer \citep{oliveirajr_etal2026_1}.
The second step to compute the RTP consists in using the predefined Neumann's layer to compute 
the transformed potential field $f_{\varepsilon}$ (Eq. \ref{eq:transformed-field-continuous-layer}) defined in the 
space domain by using $D_{t} = D_{f}$, with $D_{f} = \nicefrac{\partial^{2}}{\partial z^{2}}$, or in the Fourier domain 
by using $\alpha_{t} = \alpha_{f}$ (Eq. \ref{eq:alpha-t-generic}), with 
$\alpha_{f} = k^{2}$.
In this case, the transformed field $f_{\varepsilon}$ (Eq. \ref{eq:transformed-field-continuous-layer}) 
approximates the potential field $f$ (Eq. \ref{eq:transformed-field-space}) defined by 
$D_{f} = \nicefrac{\partial^{2}}{\partial z^{2}}$, which represents the RTP anomaly. 
As we have just shown here, for the first time, the classical approach of \cite{silva1986} to compute the RTP in the 
space domain can be formally explained in terms of a Neumann's equivalent layer. 

We can also investigate the theoretical link between the RTP in the space domain, 
via equivalent-layer technique, and in the Fourier domain, via classical approach \cite[e.g.,][p. 331]{blakely1996}. 
As we have already discussed, the transformed field $f_{\varepsilon}$ computed via \ref{eq:transformed-field-continuous-layer-Fourier} 
from a previously defined equivalent layer $p_{\varepsilon}$ can also be directly 
computed from the observed potential field $d^{o}$ (Eq. \ref{eq:observed-data-space}) on a horizontal surface
$\mathcal{S}^{\dagger}$ located above the true sources, at a constant vertical coordinate $z^{\dagger}_{0}$.
To compute the RTP, we set $\hat{\mathbf{h}} = \hat{\mathbf{v}}$ in Eq. \ref{eq:p-Fourier-dipoles}, so that 
we obtain $\mathcal{F} [ p_{\varepsilon}(\mathbf{r}'') ] = k \, \mathcal{F} \left[ U_{\varepsilon}(\mathbf{r}'') \right]$.
If we additionally consider that the surface $\mathcal{S}''$ is located above the sources, at a constant vertical coordinate
$z''_{0} < z^{t}_{0}$, where $z^{t}_{0}$ defines the top of the true sources, we obtain by
letting $k_{\varepsilon}$ goes to infinity in Eq. \ref{eq:potential-U-epsilon-Fourier-layer-1} or 
by letting $\mu_{\varepsilon}$ goes to zero in Eq. \ref{eq:potential-U-epsilon-Fourier-layer-2} that 
$\mathcal{F} \left[ U_{\varepsilon}(\mathbf{r}'') \right] = \mathcal{F} \left[ U(\mathbf{r}'') \right]$.
Finally, by noting that the product 
$\mathcal{F} \left[ U(\mathbf{r}'') \right] \mathcal{F} [ \Psi(\mathbf{r}^{\dagger}, \mathbf{r}'') ] = 
\nicefrac{\mathcal{F} \left[ U(\mathbf{r}^{\dagger}) \right]}{k}$, where 
$\mathcal{F} [ \Psi(\mathbf{r}^{\dagger}, \mathbf{r}'') ] = \nicefrac{e^{-k \, (z''_{0} - z^{\dagger}_{0})}}{k}$ 
\citep[e.g.,][p. 275]{blakely1996}, and that 
$\mathcal{F} \left[ d^{o}(\mathbf{r}^{\dagger}) \right] = \alpha_{d} \, \mathcal{F} \left[ U(\mathbf{r}^{\dagger}) \right]$,
where
$\alpha_{d} = \left(\mathbf{k}^{\top} \hat{\mathbf{u}} \right)\left(\mathbf{k}^{\top} \hat{\mathbf{v}} \right)$,
we get from Eq. \ref{eq:transformed-field-continuous-layer-Fourier} that the transformed potential field
$f_{\varepsilon}$ in the Fourier domain assumes the following form when $\alpha_{f} = k^{2}$:
\begin{equation}
	\mathcal{F} [ f_{\varepsilon}(\mathbf{r}^{\dagger}) ] =
	\tfrac{k^{2}}{\left(\mathbf{k}^{\top} \hat{\mathbf{u}} \right) \, \left(\mathbf{k}^{\top} \hat{\mathbf{v}} \right)} \, 
	\mathcal{F} \left[ d^{o}(\mathbf{r}^{\dagger}) \right] \: .
	\label{eq:transformed-field-continuous-layer-single-step-RTP}
\end{equation}
Equation \ref{eq:transformed-field-continuous-layer-single-step-RTP} represents the
classical approach to compute the RTP in the Fourier domain \citep[e.g.,][p. 331]{blakely1996}.
We stress two important aspects of this result. The first concerns the theoretical link 
between the space- and Fourier-domain approaches to compute the RTP. To the authors knowledge, this link has never been fully 
detailed in the literature.
The second important aspect of our theoretical approach is that it provides the 
reason by which the equivalent-layer technique ultimately suffers the same instability as 
the Fourier-domain RTP operator at low latitudes. In this case, the estimated equivalent layer of dipoles 
presents striations, as have already been pointed out by \cite{li_etal2014}. 
In the Fourier domain, these artifacts are directly related to the term 
$\nicefrac{\left(\mathbf{k}^{\top} \hat{\mathbf{v}} \right)}{\left(\mathbf{k}^{\top} \hat{\mathbf{h}} \right)}$ 
in equation \ref{eq:p-Fourier-dipoles}. This term produces, in the equivalent layer of dipoles, a 
numerical instability similar to that caused by the term 
$\nicefrac{k^{2}}{\left(\mathbf{k}^{\top} \hat{\mathbf{u}} \right) \, \left(\mathbf{k}^{\top} \hat{\mathbf{v}} \right)}$ in Eq. 
\ref{eq:transformed-field-continuous-layer-single-step-RTP},
which is the base of the classical approach to compute the RTP in the Fourier domain \citep[e.g.,][p. 331]{blakely1996}.

\subsection{Generalized equivalent sources}

We get from Eq. \ref{eq:p-Fourier} that the equivalent layer $p_{\varepsilon}$ obtained by using
$\alpha_{a} = 1$ is given by
\begin{equation}
	\mathcal{F} [ p_{\varepsilon}(\mathbf{r}'') ] =
	k \, \alpha_{d} \, \mathcal{F} \left[ U_{\varepsilon}(\mathbf{r}'') \right]  \: ,
	\label{eq:p-Fourier-generalized-eqs}
\end{equation}
which approximates the vertical derivative of the observed potential field on the layer plane $\mathcal{S}''$. 
Such layers produce potential-field transformations (Eq(s). \ref{eq:transformed-field-continuous-layer} and 
\ref{eq:transformed-field-continuous-layer-Fourier}) defined by $\alpha_{t} = \nicefrac{\alpha_{f}}{\alpha_{d}}$ 
(Eq. \ref{eq:alpha-t-generic}). Besides, they produce a predicted potential field $d_{\varepsilon}$ 
(Eq. \ref{eq:predicted-data-continuous-layer}) analogous to the gravitational potential.
This approach was proposed by \cite{cordell1992} and recently adapted by 
\cite{soler_uieda2021} to perform gravity data processing. 
Those authors considered only the particular case in which the observed potential field $d^{o}$ 
(Eq. \ref{eq:observed-data-space}) is gravity disturbance data, i.e., the vertical derivative 
$\tfrac{\partial}{\partial z} U$ of the potential $U$ (Eq. \ref{eq:potential-U-space}).
In this particular case, $\alpha_{d} \equiv \alpha(1, \hat{\mathbf{z}}) = k$ (Eq. \ref{eq:diff-operators-generic})
and Eq. \ref{eq:p-Fourier} leads to an equivalent layer $p_{\varepsilon}$ that
(i) approximates the second derivative $\tfrac{\partial^{2}}{\partial z^{2}} U$ of $U$ (Eq. \ref{eq:potential-U-space}) 
and
(ii) results in potential-field transformations defined by $\alpha_{t} = \nicefrac{\alpha_{f}}{k}$ (Eq. \ref{eq:alpha-t-generic}).

\subsection{Alternative approach to compute the RTP}

\cite{guspi_novara2009} proposed an alternative two-step approach to compute the RTP via equivalent-layer technique. 
They consider that the observed potential field $d^{o}$ (Eq. \ref{eq:observed-data-space}) is satisfactorily described 
by the differential operator $D_{d} \equiv D(1, \hat{\mathbf{u}}) \, D(1, \hat{\mathbf{v}})$
(Eq. \ref{eq:diff-operators-generic}), where $\hat{\mathbf{u}}$ defines the constant direction of the main geomagnetic
field at the study area and $\hat{\mathbf{v}}$ the constant total-magnetization direction of the true magnetic sources.
In this case, we get from Eq. \ref{eq:diff-operators-generic} that 
$D_{d} = \hat{\mathbf{u}}^{\top} \mathbf{H}_{U}(\mathbf{r}) \, \hat{\mathbf{v}}$,
where $\mathbf{H}_{U}$ is the Hessian matrix \citep[e.g.,][p. 158]{kaplan2002} of
the potential $U$ (Eq. \ref{eq:potential-U-space}), and 
$\alpha_{d} = \left(\mathbf{k}^{\top} \hat{\mathbf{u}} \right)\left(\mathbf{k}^{\top} \hat{\mathbf{v}} \right)$.

The first step of the alternative method proposed by \cite{guspi_novara2009} consists in obtaining an equivalent layer 
$p_{\varepsilon}$ (Eq. \ref{eq:predicted-data-continuous-layer}) by considering that
\begin{equation}
	D_{a}\Psi(\mathbf{r}, \mathbf{r}'') = \tfrac{\left( z''_{0} - z \right)}{2 \, \pi} \, 
	\ln{\Big[} \left( z''_{0} - z \right) + \| \mathbf{r}'' - \mathbf{r} \| {\Big]} - 
	\tfrac{\| \mathbf{r}'' - \mathbf{r} \|}{2 \, \pi} \: ,
	\label{eq:kernel-Fourier-Guspi}
\end{equation}
which represents the double integral of the function $\Psi$ (Eq. \ref{eq:Psi}) along the vertical direction $z$.
In this case, we have that $\alpha_{a} = \nicefrac{(\mathbf{k}^{\top} \hat{\mathbf{u}}) (\mathbf{k}^{\top} \hat{\mathbf{h}})}{k^{2}}$, 
where the term $\nicefrac{1}{k^{2}}$ represents the double integration along the vertical direction. 
Besides, Eq. \ref{eq:p-Fourier} leads to
\begin{equation}
	\mathcal{F} [ p_{\varepsilon}(\mathbf{r}'') ] =
	k^{3} \, \tfrac{\left(\mathbf{k}^{\top} \hat{\mathbf{v}} \right)}{\left(\mathbf{k}^{\top} \hat{\mathbf{h}} \right)} \,
	\mathcal{F} \left[ U_{\varepsilon}(\mathbf{r}'') \right] \: .
	\label{eq:p-Fourier-guspi}
\end{equation}
The second step of the method proposed by \cite{guspi_novara2009} consists in approximates the RTP anomaly by
the transformed potential field $f_{\varepsilon}$ obtained with $D_{t}\Psi = \Psi$ in Eq. 
\ref{eq:transformed-field-continuous-layer}, which corresponds to $\alpha_{t} = 1$ (Eq. \ref{eq:alpha-t-generic}).
In this case, by using the identity 
$\mathcal{F} \left[ U_{\varepsilon}(\mathbf{r}'') \right] \mathcal{F} [ \Psi(\mathbf{r}^{\dagger}, \mathbf{r}'') ] = 
\nicefrac{\mathcal{F} \left[ U_{\varepsilon}(\mathbf{r}^{\dagger}) \right]}{k}$, where 
$\mathcal{F} [ \Psi(\mathbf{r}^{\dagger}, \mathbf{r}'') ] = \nicefrac{e^{-k \, (z''_{0} - z^{\dagger}_{0})}}{k}$ 
\citep[e.g.,][p. 275]{blakely1996}, 
we get from Eq. \ref{eq:transformed-field-continuous-layer-Fourier}
that the transformed potential field $f_{\varepsilon}$ assumes the following form in the Fourier domain:
\begin{equation}
	\mathcal{F} [ f_{\varepsilon}(\mathbf{r}^{\dagger}) ] =
	k^{2} \, \tfrac{\left(\mathbf{k}^{\top} \hat{\mathbf{v}} \right)}{\left(\mathbf{k}^{\top} \hat{\mathbf{h}} \right)}
	\, \mathcal{F} \left[ U_{\varepsilon}(\mathbf{r}^{\dagger}) \right] \: , 
	\quad z^{\dagger}_{0} < z''_{0} \: .
	\label{eq:transformed-field-Guspi-Fourier}
\end{equation}
Equation \ref{eq:transformed-field-Guspi-Fourier} shows, in the Fourier domain, that the transformed potential field 
$\mathcal{F} [ f_{\varepsilon}(\mathbf{r}^{\dagger}) ]$ actually approximates the RTP anomaly 
(given by $k^{2} \mathcal{F} \left[ U(\mathbf{r}^{\dagger}) \right]$) for the particular case in which 
$\hat{\mathbf{h}} = \hat{\mathbf{v}}$, i.e., for the case in which the constant magnetization direction of 
the equivalent layer coincides with that of the true sources.
We have shown here that the theoretical framework proposed by \cite{oliveirajr_etal2026_1} can be used to formally 
explain the alternative method proposed by \cite{guspi_novara2009}.
